% !TEX root = ../../main.tex

\chapter{卷积网络的训练}

至此，我们已经完整搭建了卷积神经网络的架构。那么，卷积神经网络又该如何进行训练呢？实际上，CNN的训练过程与全连接网络并没有本质区别，仍然遵循“前向传播$\rightarrow$计算损失$\rightarrow$反向传播$\rightarrow$参数更新”的基本流程。不同的是，CNN中增加了卷积和池化等新的运算，因此在反向传播时，需要进一步掌握这些运算对应的梯度计算方法。下面以\cref{fig:cnn}所示的简单卷积神经网络为例，完整走一遍这一训练流程。

\begin{figure}[H]
    \boxedimage[width=0.95\textwidth]{figures/part_02/p2_ch02/cnn.pdf}
    \caption{一个简单的卷积神经网络。3通道$12\times 12$输入依次经过两组卷积与最大池化，再经64、32两个全连接层，最终输出5个类别得分。}
    \label{fig:cnn}
\end{figure}

我们首先进行前向传播。如\cref{fig:cnn}所示，该网络接收一个包含$3$个通道的$12\times12$输入。输入首先经过卷积层Conv1，使用$6$个$3\times3\times3$卷积核进行特征提取，得到$6$张$12\times12$的特征图；随后通过$2\times2$的最大池化，将每张特征图的空间尺寸压缩为$6\times6$，得到$6\times6\times6=216$个特征值。接下来，第二个卷积层Conv2使用$12$个$3\times3\times6$卷积核，步长$s=1$，填充$p=1$，进一步提取更高层次的特征，得到$12$张$6\times6$的特征图；再经过第二个$2\times2$最大池化层，将空间尺寸进一步压缩为$3\times3$，最终得到$12\times3\times3=108$个特征值。将这些特征展平后，依次送入包含$64$个和$32$个神经元的两层全连接层，最后通过包含$5$个神经元的输出层得到5个类别的logits，再通过Softmax将其转换为各类别的预测概率。为简洁起见，\cref{fig:cnn}中没有画出激活函数；本文默认每个卷积层后都接ReLU激活函数。后文反向传播中的$\Delta_k$均指已经包含$\mathrm{ReLU}'(\cdot)$影响的梯度。

\section{分类损失函数}

之后，我们进一步计算模型的损失。虽然前面介绍的均方误差（Mean Squared Error，MSE）也可以用于分类任务，但并不是理想的选择。以当前的5分类任务为例，假设一张图片的真实类别是第4类，那它的真实标签可以表示为一个长度为$5$的向量：
$$
y=[0,0,0,1,0]
$$
假设模型预测得到的概率分布为：
$$
\hat{y}=[0.02,0.03,0.01,0.90,0.04]
$$
如果使用MSE计算损失，则有：
$$
\mathcal{L}_{\mathrm{MSE}} = \frac{1}{5}\sum_{i=1}^{5}(y_i-\hat{y}_i)^2 \approx0.0026
$$
现在假设模型对第4类的预测概率从$90\%$降到$10\%$，并将减少的$80\%$概率分配给其他4个类别：
$$
\hat{y}=[0.21,0.20,0.19,0.10,0.30]
$$
此时，MSE为：
$$
\mathcal{L}_{\mathrm{MSE}} = \frac{1}{5}\sum_{i=1}^{5}(y_i-\hat{y}_i)^2 \approx0.2040
$$
而如果模型进一步地，将第4类的预测概率降到$1\%$，并将减少的$9\%$概率分配给其他类别：
$$
\hat{y}=[0.25,0.24,0.25,0.01,0.25]
$$
此时，MSE为：
$$
\mathcal{L}_{\mathrm{MSE}}=\frac{1}{5}\sum_{i=1}^{5}(y_i-\hat{y}_i)^2\approx0.2450
$$
可以看到，当模型对真实类别的预测概率从$10\%$进一步降到$1\%$时，MSE虽然还会继续增大，但增幅已经明显减小。也就是说，即使模型已经非常确信地做出了错误预测，MSE的惩罚也不够明显。因此，对于分类任务，我们需要一种对预测概率变化更加敏感的损失函数。

\textbf{交叉熵（Cross-Entropy）}就是分类任务中常用的一种损失函数。对于上面的5分类例子，真实标签$y=[0,0,0,1,0]$，模型预测概率$\hat{y}=[0.02,0.03,0.01,0.90,0.04]$。对于一个包含$K$个类别的分类任务，交叉熵可以写为：
$$
\mathcal{L} = -\sum_{i=1}^{K}y_i\log(\hat{y}_i)
$$
由于真实标签中只有第4类对应的位置为$1$，其余位置均为$0$，因此上式可以进一步化简为：
$$
\mathcal{L}=-\log(\hat{y}_4)=-\log(0.90)\approx0.105
$$
对于前面相同的三个预测结果，当真实类别为第4类时，如果模型分别给第4类分配$90\%$、$10\%$和$1\%$的概率，那么对应的交叉熵分别为：
$$
\begin{aligned}
-\log(0.9) &\approx 0.105\\
-\log(0.1) &\approx 2.303\\
-\log(0.01) &\approx 4.605
\end{aligned}
$$
可以看到，当模型给真实类别的概率从$10\%$进一步降到$1\%$时，交叉熵从约$2.30$迅速增加到约$4.61$，能够明显加大对错误且过于自信的预测的惩罚。因此，对于分类任务，交叉熵通常比MSE更加合适。

\section{反向传播}

接下来，我们进行反向传播，计算网络中各个参数对损失的梯度。对于全连接层和激活函数，我们之前已经介绍过其反向传播过程，因此这里主要关注卷积层和池化层。

\subsection{池化层的反向传播}

对于池化层，前向传播时没有需要学习的参数，因此反向传播的任务不是更新参数，而是将梯度继续传回前一层。\cref{fig:cnn}中，第二层池化层Pool2输出的特征图大小为$12\times3\times3$，经过Flatten后得到$12\times3\times3=108$个特征值，再送入包含$64$个神经元的全连接层。反向传播时，梯度首先从全连接层传回Flatten层。由于Flatten只改变数据的排列方式，并不进行实际的数值计算，因此这$108$个梯度会重新对应到原来的$12\times3\times3$特征图。于是，对于Pool2而言，我们需要处理的就是这$12\times3\times3$个位置收到的梯度，并将它们继续传回Pool2的输入，也就是Conv2输出的$12\times6\times6$特征图。

以其中一个$2\times2$池化窗口为例，假设前向传播时该窗口的输入为：
$$
\begin{pmatrix}
1&5\\
3&2
\end{pmatrix}
$$
采用最大池化后，该窗口只保留最大值$5$，因此对应的池化输出为$5$。假设反向传播时，这个输出位置从后一层接收到的梯度为$\delta$，那么该梯度只会传回前向传播时被选中的最大值位置，而其余三个位置由于在前向传播时被舍弃而没有参与输出，因此不会接收到梯度，梯度均为$0$。于是，回传的梯度为：
$$
\begin{pmatrix}
0&\delta\\
0&0
\end{pmatrix}
$$
也就是说，最大池化在反向传播时只沿着前向传播时的“最大值路径”传递梯度。

\begin{figure}[H]
    \boxedimage[width=0.75\textwidth]{figures/part_02/p2_ch02/bp_pooling.pdf}
    \caption{最大池化的反向传播。前向传播将$12\times6\times6$的特征图通过$2\times2$、步长为$2$的最大池化得到$12\times3\times3$的特征图；反向传播时，$12\times3\times3$的梯度根据前向传播时记录的最大值位置重新分配，恢复为$12\times6\times6$的梯度。}
    \label{fig:bp_pooling}
\end{figure}

如\cref{fig:bp_pooling}所示，这一过程会应用到整张特征图的所有池化窗口：$12\times3\times3$的梯度会按照前向传播时记录的最大值位置，分别传回对应的$2\times2$区域，最终得到$12\times6\times6$的梯度。由于每个通道都独立进行池化，因此这一过程会分别应用于$12$个通道。

对于平均池化，处理方式则有所不同。由于前向传播时计算的是窗口内所有元素的平均值，因此反向传播时，传回的梯度会平均分配给窗口中的每一个位置。例如，一个$2\times2$窗口经过平均池化得到
$$
y=\frac{x_1+x_2+x_3+x_4}{4}
$$
如果从后一层传回的梯度为$\delta$，那么根据链式法则，窗口中每个元素接收到的梯度都是$\frac{\delta}{4}$，因此回传的梯度矩阵为：
$$
\begin{pmatrix}
\frac{\delta}{4}&\frac{\delta}{4}\\
\frac{\delta}{4}&\frac{\delta}{4}
\end{pmatrix}
$$

\subsection{卷积层的反向传播}

对于卷积层，情况与池化层有所不同。卷积层包含需要学习的权重和偏置，因此反向传播时不仅需要将梯度继续传回前一层，还需要计算损失函数对卷积核权重和偏置的梯度。这里以\cref{fig:cnn}中的Conv2为例。按照前面的网络结构，Conv2接收来自Pool1的$6\times6\times6$特征图，使用$12$个$3\times3\times6$的卷积核，步长$s=1$、填充$p=1$，最终得到$12\times6\times6$的输出特征图。如\cref{fig:bp_convolution}所示，在反向传播过程中，来自下游Pool2的梯度会同时用于三个方向的计算：一方面计算卷积核的权重梯度，另一方面计算偏置梯度，同时还需要计算损失函数对Conv2输入的梯度，并将其继续传递给前面的Pool1。

\begin{figure}[H]
    \boxedimage[width=0.75\textwidth]{figures/part_02/p2_ch02/bp_convolution.pdf}
    \caption{卷积层Conv2的反向传播。来自Pool2的梯度同时用于计算卷积核的权重梯度、偏置梯度以及输入梯度：将梯度与对应输入逐位置相乘并累加，得到权重梯度；对空间位置上的梯度求和，得到偏置梯度；将梯度与对应卷积核权重进行加权并累加，得到输入梯度，并继续传回Pool1。}
    \label{fig:bp_convolution}
\end{figure}

为了便于说明，我们先只关注Conv2中的一个输出通道。记第$k$个卷积核为$W_k$，由于输入包含$6$个通道，因此$W_k$实际上由$6$个$3\times3$的权重矩阵组成：
$$
W_k=\left(W_{k,1},W_{k,2},\ldots,W_{k,6}\right)
$$
其中，$W_{k,c}$表示第$k$个卷积核在第$c$个输入通道上的$3\times3$权重矩阵。为了进一步观察其中一个权重的梯度，我们固定一个输入通道$c$，将对应的权重矩阵记为：
$$
W_{k,c}=\begin{pmatrix}
w_{k,c,11}&w_{k,c,12}&w_{k,c,13}\\
w_{k,c,21}&w_{k,c,22}&w_{k,c,23}\\
w_{k,c,31}&w_{k,c,32}&w_{k,c,33}
\end{pmatrix}
$$

在反向传播时，我们首先需要得到损失函数对Conv2线性输出的梯度。记第$k$个输出通道在位置$(i,j)$处的线性输出为$Z_k(i,j)$，则有：
$$
\Delta_k(i,j)=\frac{\partial\mathcal{L}}{\partial Z_k(i,j)}
$$
这里的$\Delta_k$已经包含ReLU的反向传播，即已经乘上$\mathrm{ReLU}'(\cdot)$。后文使用$\Delta_k$进行卷积层梯度计算时，不再单独展开ReLU的梯度。Pool2在前向传播时将$12\times6\times6$的特征图通过$2\times2$最大池化得到$12\times3\times3$的输出，因此反向传播时，Pool2会将$12\times3\times3$的梯度重新分配到$12\times6\times6$的位置。假设对于第$k$个输出通道，传回Conv2的梯度矩阵为：
$$
\Delta_k=\left(\begin{array}{cc|cc|cc}
0&\delta_{11}&0&0&0&\delta_{13}\\
0&0&\delta_{12}&0&0&0\\ \hline
\delta_{21}&0&0&\delta_{22}&0&0\\
0&0&0&0&\delta_{23}&0\\ \hline
0&\delta_{31}&\delta_{32}&0&0&0\\
0&0&0&0&0&\delta_{33}
\end{array}
\right)
$$
其中，$\delta_{11},\ldots,\delta_{33}$分别来自Pool2输出的$3\times3$个位置。

接下来，我们以$W_{k,c}$中左上角的权重$w_{k,c,11}$为例。当第$k$个卷积核位于输出位置$(i,j)$时，该权重会与第$c$个输入通道中当前$3\times3$感受野左上角的元素$\tilde{X}_c(i,j)$相乘，因此该位置对$w_{k,c,11}$产生的梯度为：
$$
\Delta_k(i,j)\tilde{X}_c(i,j)
$$
由于同一个权重会随着卷积核在特征图上移动，在全部$6\times6=36$个位置重复使用，因此所有位置产生的梯度都需要累加：
$$
\frac{\partial\mathcal{L}}{\partial w_{k,c,11}}=\sum_{i=1}^{6}\sum_{j=1}^{6}\Delta_k(i,j)\tilde{X}_c(i,j)
$$
代入上面的$\Delta_k$矩阵，只有其中$9$个位置具有非零梯度，因此可以进一步写成：
$$
\begin{aligned}
\frac{\partial\mathcal{L}}{\partial w_{k,c,11}}
={}&
\delta_{11}\tilde{X}_c(1,2)+\delta_{12}\tilde{X}_c(2,3)+\delta_{13}\tilde{X}_c(1,6)+\delta_{21}\tilde{X}_c(3,1)+\delta_{22}\tilde{X}_c(3,4)\\
&+\delta_{23}\tilde{X}_c(4,5)+\delta_{31}\tilde{X}_c(5,2)+\delta_{32}\tilde{X}_c(5,3)+\delta_{33}\tilde{X}_c(6,6)
\end{aligned}
$$
这里的$\tilde{X}_c$表示经过零填充后的第$c$个输入通道。这里$X_c$的尺寸为$6\times6$，经过$p=1$的零填充后，$\tilde X_c$的尺寸变为$8\times8$。而由于卷积使用了零填充，靠近输入边界的位置可能对应填充区域，因此相应位置的$\tilde{X}_c$值可能为$0$。

对于卷积核中的其他权重，计算过程完全相同，只是它们在每个空间位置对应的输入元素不同。例如，权重$w_{k,c,12}$对应当前$3\times3$感受野第一行第二列的输入值，因此有：
$$
\frac{\partial\mathcal{L}}{\partial w_{k,c,12}}=\sum_{i=1}^{6}\sum_{j=1}^{6}\Delta_k(i,j)\tilde{X}_c(i,j+1)
$$
同理，可以得到$W_{k,c}$中其余权重的梯度。因此，对于第$k$个输出通道和第$c$个输入通道，对应的$3\times3$权重梯度矩阵可以写成：
$$
\frac{\partial\mathcal{L}}{\partial W_{k,c}}=
\begin{pmatrix}
\frac{\partial\mathcal{L}}{\partial w_{k,c,11}}&
\frac{\partial\mathcal{L}}{\partial w_{k,c,12}}&
\frac{\partial\mathcal{L}}{\partial w_{k,c,13}}\\
\frac{\partial\mathcal{L}}{\partial w_{k,c,21}}&
\frac{\partial\mathcal{L}}{\partial w_{k,c,22}}&
\frac{\partial\mathcal{L}}{\partial w_{k,c,23}}\\
\frac{\partial\mathcal{L}}{\partial w_{k,c,31}}&
\frac{\partial\mathcal{L}}{\partial w_{k,c,32}}&
\frac{\partial\mathcal{L}}{\partial w_{k,c,33}}
\end{pmatrix}
$$

由于第$k$个卷积核需要同时覆盖输入的$6$个通道，因此还需要分别计算这$6$个输入通道对应的权重梯度。最终，第$k$个卷积核的完整梯度为：
$$
\frac{\partial\mathcal{L}}{\partial W_k}\in\mathbb{R}^{6\times3\times3}
$$
Conv2共有$12$个卷积核，因此对这$12$个卷积核分别计算后，可以得到整个卷积层的权重梯度：
$$
\frac{\partial\mathcal{L}}{\partial W_2}\in\mathbb{R}^{12\times6\times3\times3}.
$$

卷积核的偏置同样会在整个$6\times6$的输出特征图上重复使用。对于第$k$个输出通道，偏置$b_k$会参与该通道每一个空间位置的计算，因此其梯度也需要将所有位置的梯度累加：
$$
\frac{\partial\mathcal{L}}{\partial b_k}=\sum_{i=1}^{6}\sum_{j=1}^{6}\Delta_k(i,j)
$$
根据上面的$\Delta_k$矩阵，可以得到：
$$
\frac{\partial\mathcal{L}}{\partial b_k}=\delta_{11}+\delta_{12}+\delta_{13}+\delta_{21}+\delta_{22}+\delta_{23}+\delta_{31}+\delta_{32}+\delta_{33}
$$
对$12$个输出通道分别计算后，可以得到整个Conv2偏置的梯度：
$$
\frac{\partial\mathcal{L}}{\partial b_2}\in\mathbb{R}^{12}
$$

到这里，我们已经计算出了Conv2自身参数$W_2$和$b_2$的梯度。但反向传播还没有结束，因为Conv2的输入来自前面的Pool1，所以还需要进一步计算损失函数对Conv2输入的梯度，并将其继续传回Pool1。

与计算卷积核参数梯度时类似，一个输入位置通常会参与多个卷积输出位置的计算，因此会同时从多个输出位置接收到梯度。同时，由于Conv2的输入包含$6$个通道，每个输入通道中的一个位置都会参与全部$12$个输出通道的计算，因此来自不同卷积核的梯度还需要进一步累加。为了说明这一过程，我们固定Pool1输出的第$c$个通道。需要注意的是，Pool1的完整输出是一个$6\times6\times6$的特征图，其中有$6$个$6\times6$的通道；固定第$c$个通道后，我们实际考察的是其中一张$6\times6$的特征图。现在考虑其中的一个内部位置$X_c(3,3)$。由于Conv2使用$3\times3$卷积核并采用$p=1$的零填充，因此$X_c(3,3)$会参与输出位置$(2,2)$到$(4,4)$对应的$9$个卷积计算。

\begin{figure}[H]
    \boxedimage[width=0.6\textwidth]{figures/part_02/p2_ch02/bp_input.pdf}
    \caption{输入位置$X_c(3,3)$参与的9个卷积计算。红框为各输出位置（$(2,2)$--$(4,4)$）的感受野，粉红格为$(3,3)$，标题$w_{k,c,ij}$为对应的核权重；该位置的梯度由9路梯度按这些权重累加。}
    \label{fig:bp_input}
\end{figure}

如\cref{fig:bp_input}所示，这$9$个输出位置分别对应卷积核中的不同权重。对于第$k$个卷积核，它们之间的具体对应关系如下：

\begin{table}[H]
    \centering
    \caption{$X_c(3,3)$在不同卷积输出位置对应的卷积核权重}
    \label{tab:conv2_input_gradient}
    \begin{tabular}{cc}
        \toprule
        \thead{输出位置} & \thead{$X_c(3,3)$对应的权重} \\
        \midrule
        $(2,2)$ & $w_{k,c,33}$ \\
        $(2,3)$ & $w_{k,c,32}$ \\
        $(2,4)$ & $w_{k,c,31}$ \\
        $(3,2)$ & $w_{k,c,23}$ \\
        $(3,3)$ & $w_{k,c,22}$ \\
        $(3,4)$ & $w_{k,c,21}$ \\
        $(4,2)$ & $w_{k,c,13}$ \\
        $(4,3)$ & $w_{k,c,12}$ \\
        $(4,4)$ & $w_{k,c,11}$ \\
        \bottomrule
    \end{tabular}
\end{table}

因此，第$k$个卷积核从$X_c(3,3)$接收到的梯度为：
\begin{align*}
\frac{\partial\mathcal{L}}{\partial X_c(3,3)}\bigg|_k
={}&
\Delta_k(2,2)w_{k,c,33}+\Delta_k(2,3)w_{k,c,32}+\Delta_k(2,4)w_{k,c,31}\\
&+\Delta_k(3,2)w_{k,c,23}+\Delta_k(3,3)w_{k,c,22}+\Delta_k(3,4)w_{k,c,21}\\
&+\Delta_k(4,2)w_{k,c,13}+\Delta_k(4,3)w_{k,c,12}+\Delta_k(4,4)w_{k,c,11}
\end{align*}
由于$X_c(3,3)$同时参与Conv2的$12$个输出通道的计算，因此还需要将这$12$个卷积核产生的梯度全部累加：
$$
\frac{\partial\mathcal{L}}{\partial X_c(3,3)}
=
\sum_{k=1}^{12}
\frac{\partial\mathcal{L}}{\partial X_c(3,3)}
\bigg|_k.
$$
按照相同的方法，可以计算出Pool1输出特征图中所有位置的梯度，最终得到一个与Conv2输入尺寸相同的$6\times6\times6$梯度张量。这个梯度随后传回Pool1，再由Pool1按照最大池化的反向传播规则，将梯度继续传给Conv1。Conv1收到梯度后，可以采用相同的思路计算自身的权重和偏置梯度，并进一步将梯度传向网络输入。因此，对于当前这个包含两组卷积和池化的网络，梯度会沿着下面的路径逐层反向传播：
$$
\mathcal{L}\rightarrow\mathrm{FC}\rightarrow\mathrm{Pool2}\rightarrow\mathrm{Conv2}\rightarrow\mathrm{Pool1}\rightarrow\mathrm{Conv1}
$$

当网络中各层的梯度都计算完成后，所有可学习参数的梯度也就得到了。接下来进入参数更新阶段，这一步与前面介绍的全连接网络完全相同。对于任意可学习参数$\theta$，最基本的梯度下降更新公式为：
$$
\theta\leftarrow\theta-\eta\nabla_{\theta}\mathcal{L}
$$
可以看到，CNN与全连接网络在参数更新方式上并没有本质区别，真正不同的是梯度的计算方式。在CNN中，最大池化层根据前向传播时记录的最大值位置传递梯度，卷积层则利用这些梯度计算共享权重和偏置的梯度，并继续将梯度传向前面的网络层。完成梯度计算后，仍然可以使用SGD、动量法或Adam等优化算法更新网络参数。

\section{卷积的统一理解}

前面我们已经从网络结构和训练过程两个角度学习了卷积神经网络：卷积层通过局部连接和权值共享提取空间特征，池化层进一步压缩特征图，反向传播则负责计算各层参数的梯度。到这里，有必要再从一个更统一的角度看待卷积。卷积并不是一种完全不同于全连接层的计算方式，它本质上仍然是在进行加权求和，只是对连接方式和参数使用方式增加了结构约束。理解这一点，可以把前面介绍的卷积计算、参数量减少、平移性质以及反向传播统一起来。

\subsection{结构化线性变换}

我们再回顾一下全连接层。对于输入向量$x$，全连接层通常写成：
$$
z=Wx+b
$$
其中$W$是一个稠密的权重矩阵，每个输出神经元都可以连接到输入的所有维度。因此，如果输入包含$N_{\mathrm{in}}$个特征，输出包含$N_{\mathrm{out}}$个神经元，那么权重矩阵$W$就包含$N_{\mathrm{in}} \cdot N_{\mathrm{out}}$个参数。也就是说，全连接层允许任意输入特征与任意输出特征建立连接，网络并不对连接关系施加额外限制。

卷积层的计算形式在这一点上并没有本质区别。以Conv2为例，每个输出位置都会对输入中的一个局部区域进行加权求和。假设输入为$X$，卷积核为$W_k$，那么第$k$个输出通道在位置$(i,j)$处的线性输出，本质上仍然可以写成：
$$
Z_k(i,j)=\sum_c\sum_r\sum_s w_{k,c,rs}\tilde X_c(i+r-1,j+s-1)+b_k
$$
它依然是在做“输入$\times$权重，再求和，加上偏置”，与全连接层的计算思想完全一致。不同之处在于，卷积层并不允许一个输出位置连接到所有输入位置，而只允许它连接到一个局部感受野；与此同时，不同空间位置使用的是同一组卷积核参数。

为了更直观地理解这一点，我们先看一个最简单的二维卷积例子。假设输入是一张$3\times3$的单通道图像：
$$
X=\begin{pmatrix}
x_{11}&x_{12}&x_{13}\\
x_{21}&x_{22}&x_{23}\\
x_{31}&x_{32}&x_{33}
\end{pmatrix}
$$
使用一个$2\times2$的卷积核：
$$
K=\begin{pmatrix}
w_{11}&w_{12}\\
w_{21}&w_{22}
\end{pmatrix}
$$
步长$s=1$，不使用填充$p=0$。此时，卷积核在输入上滑动，可以得到一个$2\times2$的输出特征图：
$$
Z=\begin{pmatrix}
z_{11}&z_{12}\\
z_{21}&z_{22}
\end{pmatrix}
$$
到这里，我们看到的仍然是熟悉的“卷积核在输入上滑动”。

接下来，我们换一种表示方式。将上述$3\times3$输入按照从左到右、从上到下的顺序展平，可以得到
$$
x=\begin{pmatrix}
x_{11}\\
x_{12}\\
x_{13}\\
x_{21}\\
x_{22}\\
x_{23}\\
x_{31}\\
x_{32}\\
x_{33}
\end{pmatrix}
$$
同样，将$2\times2$的输出展平，可以得到：
$$
z=\begin{pmatrix}
z_{11}\\
z_{12}\\
z_{21}\\
z_{22}
\end{pmatrix}
$$
于是，原来的卷积操作就可以重新写成一个矩阵乘法。根据前面四个输出位置的计算公式，有
$$
\begin{pmatrix}
z_{11}\\
z_{12}\\
z_{21}\\
z_{22}
\end{pmatrix}
=
\begin{pmatrix}
w_{11}&w_{12}&0&w_{21}&w_{22}&0&0&0&0\\
0&w_{11}&w_{12}&0&w_{21}&w_{22}&0&0&0\\
0&0&0&w_{11}&w_{12}&0&w_{21}&w_{22}&0\\
0&0&0&0&w_{11}&w_{12}&0&w_{21}&w_{22}
\end{pmatrix}
\begin{pmatrix}
x_{11}\\
x_{12}\\
x_{13}\\
x_{21}\\
x_{22}\\
x_{23}\\
x_{31}\\
x_{32}\\
x_{33}
\end{pmatrix}
$$

此时我们就可以清楚地看到，卷积同样可以表示成
$$
z=Wx+b,
$$
只不过这里的$W$并不是普通的稠密矩阵，而是一个具有特殊结构的权重矩阵。在这个例子中，$W$是一个$4\times9$的矩阵。由于每个输出位置只连接输入中的一个$2\times2$局部区域，因此大量位置没有实际连接，对应的权重就是$0$。例如，$z_{11}$只与$x_{11},x_{12},x_{21},x_{22}$有关，而不会直接连接到$x_{13},x_{23},x_{31},x_{32},x_{33}$。这就是卷积中的局部连接。

另一方面，可以观察到$w_{11}$、$w_{12}$、$w_{21}$和$w_{22}$会在矩阵$W$的不同位置反复出现。例如，第一个输出位置和第二个输出位置都使用相同的$w_{11}$，只是它们分别作用于不同的输入位置。这些位置上的$w_{11}$并不是彼此独立的参数，而是同一个卷积核参数在不同空间位置被重复使用。这就是权值共享。

如果换成普通的全连接层，同样的$9$个输入和$4$个输出之间理论上可以建立完全连接，此时权重矩阵中的$4\times9=36$个位置都可以拥有独立的参数。而在卷积层中，这些参数被局部连接和权值共享两条规则约束：一方面，大量连接被直接去除，对应位置固定为$0$；另一方面，剩余连接并不拥有彼此独立的参数，而是由同一个卷积核在不同位置重复使用。因此，虽然卷积层在展开后仍然可以写成$z=Wx+b$，但这里的$W$已经不再是普通全连接层中的任意稠密矩阵，而是一个具有稀疏性和参数共享结构的矩阵。

从这个角度看，卷积层可以理解为一个带结构约束的全连接层。全连接层允许所有输入与输出自由连接，而卷积层通过局部连接和权值共享，对这些连接施加了明确的结构约束。正是这两条约束，使卷积层在大幅减少参数的同时，保留了图像中非常重要的空间结构。更重要的是，这种结构并不只是为了减少参数，而是在利用图像本身的先验：局部区域中的像素具有更强的相关性，而相同的局部模式又可能出现在图像的不同位置，因此没有必要为每一个位置分别学习一套参数。

\subsection{矩阵视角}

如果卷积本质上仍然是一种加权求和，那么它能否像全连接层一样，写成矩阵乘法的形式呢？答案是可以的，这就是\textbf{im2col}（image to columns）的基本思想。它并没有改变卷积的计算结果，而是把原本分散在不同空间位置上的局部卷积操作重新排列成矩阵乘法，从而让卷积能够利用成熟的矩阵乘法实现。

为了更直观地理解这一过程，我们仍然使用前面的$3\times3$输入和$2\times2$卷积核。输入和卷积核分别为：
$$
X=\begin{pmatrix}
x_{11}&x_{12}&x_{13}\\
x_{21}&x_{22}&x_{23}\\
x_{31}&x_{32}&x_{33}
\end{pmatrix}, \qquad
K=\begin{pmatrix}
w_{11}&w_{12}\\
w_{21}&w_{22}
\end{pmatrix}
$$
步长$s=1$，填充$p=0$。我们知道，这个卷积会得到$2\times2$的输出，因此一共有4个空间位置。对于每一个输出位置，都可以从输入中取出一个$2\times2$的局部区域：
$$
\begin{pmatrix}
x_{11}&x_{12}\\
x_{21}&x_{22}
\end{pmatrix},
\qquad
\begin{pmatrix}
x_{12}&x_{13}\\
x_{22}&x_{23}
\end{pmatrix},
\qquad
\begin{pmatrix}
x_{21}&x_{22}\\
x_{31}&x_{32}
\end{pmatrix},
\qquad
\begin{pmatrix}
x_{22}&x_{23}\\
x_{32}&x_{33}
\end{pmatrix}
$$
每个局部区域都包含4个元素。接下来，将每个局部区域按照相同的顺序展平为一个列向量，并将4个列向量依次排列，就可以得到：
$$
X_{\mathrm{col}}=\begin{pmatrix}
x_{11}&x_{12}&x_{21}&x_{22}\\
x_{12}&x_{13}&x_{22}&x_{23}\\
x_{21}&x_{22}&x_{31}&x_{32}\\
x_{22}&x_{23}&x_{32}&x_{33}
\end{pmatrix}
$$
这里$X_{\mathrm{col}}$是一个$4\times4$的矩阵，每一列对应输入中的一个局部感受野。例如，第一列：
$$
\begin{pmatrix}
x_{11}\\
x_{12}\\
x_{21}\\
x_{22}
\end{pmatrix}
$$
对应左上角的$2\times2$区域，第二列则对应向右移动后的局部区域。可以看到，同一个输入元素可能会在$X_{\mathrm{col}}$中出现多次，例如$x_{22}$就出现在了多个局部区域中。这是因为不同的感受野之间存在重叠。

与此同时，卷积核也可以展平成一个行向量：
$$
W_{\mathrm{col}}=\begin{pmatrix}w_{11}&w_{12}&w_{21}&w_{22}\end{pmatrix}
$$
于是，原来的4次局部卷积就可以一次性写成矩阵乘法：
$$
Z_{\mathrm{col}}=W_{\mathrm{col}}X_{\mathrm{col}}
$$
即：
$$
\begin{pmatrix}
z_{11}&z_{12}&z_{21}&z_{22}
\end{pmatrix}
=
\begin{pmatrix}
w_{11}&w_{12}&w_{21}&w_{22}
\end{pmatrix}
\begin{pmatrix}
x_{11}&x_{12}&x_{21}&x_{22}\\
x_{12}&x_{13}&x_{22}&x_{23}\\
x_{21}&x_{22}&x_{31}&x_{32}\\
x_{22}&x_{23}&x_{32}&x_{33}
\end{pmatrix}.
$$
将矩阵乘法展开，可以得到：
$$
\begin{aligned}
z_{11}&=w_{11}x_{11}+w_{12}x_{12}+w_{21}x_{21}+w_{22}x_{22}\\
z_{12}&=w_{11}x_{12}+w_{12}x_{13}+w_{21}x_{22}+w_{22}x_{23}\\
z_{21}&=w_{11}x_{21}+w_{12}x_{22}+w_{21}x_{31}+w_{22}x_{32}\\
z_{22}&=w_{11}x_{22}+w_{12}x_{23}+w_{21}x_{32}+w_{22}x_{33}
\end{aligned}
$$

这与前面的卷积计算完全一致。也就是说，im2col并没有改变卷积本身，只是先把每个空间位置对应的局部区域整理出来，再将原本需要“滑动卷积核”完成的计算转换成了一次矩阵乘法。

这个例子中只有一个卷积核，因此$W_{\mathrm{col}}$只有一行。如果有多个卷积核，只需要将每个卷积核分别展平成一行，再将它们组合起来即可。例如，假设有$12$个卷积核，那么：
$$
W_{\mathrm{col}}\in\mathbb{R}^{12\times4} \qquad X_{\mathrm{col}}\in\mathbb{R}^{4\times4}
$$
于是矩阵乘法得到：
$$
Z_{\mathrm{col}}=W_{\mathrm{col}}X_{\mathrm{col}}\in\mathbb{R}^{12\times4}
$$
其中每一行对应一个卷积核产生的输出，每一列对应一个空间位置。最后，再将每一行按照原来的空间位置重新排列，就可以恢复为对应的输出特征图。

通过这个例子可以看到，im2col的核心就是将每个空间位置对应的局部感受野提取出来并排列成矩阵，再将卷积核展开后通过矩阵乘法一次性完成所有位置的卷积计算，即：
$$
Z=W_{\mathrm{col}}X_{\mathrm{col}}+b
$$
这种表示并没有改变卷积层本身的计算方式，而是将原本具有空间结构的局部运算重新组织成标准的矩阵乘法，使卷积能够利用高度优化的线性代数计算库进行高效计算。因此，卷积虽然具有局部连接和权值共享等特殊结构，但其底层计算仍然可以统一归结为熟悉的线性代数运算。

\subsection{平移等变与平移不变}

卷积之所以特别适合处理图像，不仅因为它通过局部连接和权值共享减少了参数，还因为这些结构使卷积天然具有一定的空间性质。其中，一个非常重要的性质就是\textbf{平移等变性（translation equivariance）}。简单来说，如果输入图像中的某个局部模式发生平移，那么卷积产生的特征响应也会跟着发生相应的平移，而不会因为位置改变就完全失去原来的响应模式。

这一性质可以通过\cref{fig:translation_equivariance}直观理解。图中第一行给出了一幅包含数字$7$的输入图像，卷积核则用于检测其中的横线结构。当卷积核在图像上滑动时，可以在数字$7$对应的横线位置产生较强的响应。随后，将输入中的数字$7$向左移动$2$个像素、向下移动$1$个像素，再使用完全相同的卷积核进行检测，可以看到对应的横线响应也随之发生了相同的位移。也就是说，卷积核并不需要因为特征的位置发生变化而重新学习，而是可以利用同一组参数在不同位置检测相同的局部模式。

\begin{figure}[H]
    \boxedimage[width=0.8\textwidth]{figures/part_02/p2_ch02/translation_equivariance.pdf}
    \caption{平移等变性示意图。第一行的数字$7$包含明显的横线结构，卷积核用于检测该结构；将数字$7$向左移动$2$个像素、向下移动$1$个像素后，卷积核仍能检测到对应的横线响应，且响应位置随输入同步移动。}
    \label{fig:translation_equivariance}
\end{figure}

用数学形式表示，设$T_\delta$表示将输入沿某个方向平移$\delta$，$\operatorname{Conv}$表示卷积操作，那么在步长为$1$且忽略边界效应的理想情况下，有：
$$
\operatorname{Conv}(T_\delta X)=T_\delta\bigl(\operatorname{Conv}(X)\bigr)
$$

这就是平移等变性。这里的“等变”可以简单理解为：输入怎么移动，输出就以相同方式移动。需要注意，它并不是说卷积输出完全不受位置影响，相反，卷积会保留这种位置变化，只是这种变化以对应的方式反映到了输出特征图中。这种性质与卷积核的权值共享密切相关。因为同一个卷积核会在不同空间位置重复使用，所以它学习到的是一种可以在不同位置重复检测的局部模式，而不是某个固定位置上的专属特征。例如，一个卷积核可能学习到某种边缘、纹理或者局部形状，那么这种模式无论出现在图像的左侧、中间还是右侧，只要局部结构相似，卷积核都可以产生相似的响应。正因为如此，卷积不需要为每个空间位置分别学习一套参数，就能够在整幅图像中寻找相同的局部模式。

不过，平移等变意味着特征的位置变化会被保留下来。对于图像分类任务而言，我们通常更关心某个特征是否出现，而不希望模型过度依赖它的精确位置。此时，池化层可以进一步减弱这种位置变化的影响。以$2\times2$最大池化为例，假设一个池化窗口中的特征值为：
$$
\begin{pmatrix}
1&3\\
2&1
\end{pmatrix}
$$
最大池化的结果为：
$$
\max(1,3,2,1)=3
$$
现在，如果这个较大的响应在同一个池化窗口内发生了位置变化，例如变成：
$$
\begin{pmatrix}
1&1\\
3&2
\end{pmatrix}
$$
最大池化的结果仍然是：
$$
\max(1,1,3,2)=3
$$
也就是说，虽然最大值的位置发生了变化，但只要它仍然位于同一个池化窗口中，池化后的输出就可能保持不变。

因此，池化在一定程度上引入了\textbf{局部平移不变性（local translation invariance）}。这里的“不变”与前面的“等变”正好形成对比：卷积层中，输入发生平移后，输出响应也会随之平移；而经过池化后，在一定范围内的小位置变化可能不会改变最终输出。需要注意的是，这种不变性是局部的，并不是说任意大小的平移都会产生完全相同的输出。如果一个特征从一个池化窗口移动到了另一个池化窗口，池化结果仍然可能发生变化。

从这个角度看，卷积和池化承担了不同的作用。卷积通过局部连接和权值共享，在不同位置检测相同的局部模式，同时保留这些模式的空间位置，因此体现了平移等变性；池化则通过对局部区域进行压缩，减弱特征在小范围内发生位置变化所带来的影响，从而增强局部平移不变性。因此，卷积神经网络处理图像时，并不是简单地对像素进行重复计算，而是在利用图像本身的空间结构：卷积负责发现局部模式并保留其位置信息，池化则进一步降低模型对局部精确位置的敏感程度。这种“等变与不变”的配合，使网络既能够识别图像中的局部特征，又能够在一定程度上适应这些特征在图像中的位置变化。

\subsection{反向传播再理解与过渡}

有了前面的矩阵视角，再回头看卷积层的反向传播，就会发现它与全连接层并没有本质区别。前面我们将卷积写成：
$$
Z=W_{\mathrm{col}}X_{\mathrm{col}}+b
$$
这与全连接层
$$
Z=WX+b
$$
在形式上完全一致。因此，只要应用矩阵乘法的链式法则，就可以直接得到卷积层的梯度。记：
$$
\Delta=\frac{\partial\mathcal{L}}{\partial Z}
$$
则卷积核展开后对应的权重梯度为：
$$
\frac{\partial\mathcal{L}}{\partial W_{\mathrm{col}}}=\Delta X_{\mathrm{col}}^{T}
$$
对展开后输入的梯度为：
$$
\frac{\partial\mathcal{L}}{\partial X_{\mathrm{col}}}=W_{\mathrm{col}}^{T}\Delta
$$
偏置在所有空间位置共享，因此其梯度就是各个空间位置上的梯度之和：
$$
\frac{\partial\mathcal{L}}{\partial b}=\Delta\mathbf{1} \qquad \mathbf{1}\in\mathbb{R}^{36}
$$
其中$\mathbf{1}$为全1向量，用于对$36$个空间位置上的梯度进行求和。

这三个公式与前面手动推导卷积层反向传播得到的结果完全对应。对于权重梯度，$\Delta X_{\mathrm{col}}^T$实际上就是把卷积核在不同空间位置产生的梯度贡献全部累加起来。前面我们针对某一个具体参数得到：
$$
\frac{\partial\mathcal{L}}{\partial w_{k,c,rs}}=\sum_{i=1}^{6}\sum_{j=1}^{6}\Delta_k(i,j)\tilde X_c(i+r-1,j+s-1)
$$
这里的“逐位置相乘，再求和”，正是矩阵乘法$\Delta X_{\mathrm{col}}^T$在某一个卷积核参数上的具体展开。换句话说，前面的手工推导是在逐个参数计算，而矩阵形式则一次性完成了所有卷积核参数的梯度计算。

输入梯度也是同样的道理。$W_{\mathrm{col}}^T\Delta$表示每个输出位置产生的梯度按照对应的卷积核权重传回输入。由于相邻的感受野彼此重叠，同一个输入位置通常会参与多个卷积计算，因此它会同时接收到来自多个输出位置的梯度，这些贡献最终需要累加。例如前面计算$X_c(3,3)$时，它会参与多个输出位置的卷积计算，因此其梯度需要由多路梯度按照对应的卷积核权重加权后共同累加。这正是卷积反向传播中“一个输入位置为什么会收到多路梯度”的原因。从矩阵角度看，这一过程就是$W_{\mathrm{col}}^T\Delta$；在实际恢复张量形状时，还需要将展开后的梯度按照各局部感受野对应的位置累加回原始输入空间，即完成相应的\textbf{col2im}操作。

这里还需要特别注意卷积层中的权值共享。同一个卷积核参数会在不同空间位置被重复使用，因此这些位置产生的梯度不能彼此独立，而必须全部累加，才能得到该参数最终的梯度。前面公式中的求和，以及矩阵形式中的$\Delta X_{\mathrm{col}}^T$，本质上描述的是同一件事情。也就是说，卷积层反向传播看起来比全连接层复杂，并不是因为它使用了新的求导规则，而是因为同一个参数被重复使用，使得不同空间位置产生的梯度需要汇总到同一个参数上。

从这个角度重新理解整个训练过程，就会发现卷积层与全连接层之间的联系更加清晰。前向传播时，卷积利用局部连接和权值共享，对不同空间位置的局部区域进行加权求和；反向传播时，则将各个空间位置产生的梯度按照相同的结构汇总到共享参数，并将梯度继续传向前面的输入。完成这些梯度计算后，参数更新仍然使用前面介绍过的梯度下降、SGD、动量法或Adam等方法。因此，CNN与全连接网络在参数更新规则上并没有本质区别，真正不同的是其前向计算的结构，以及由此决定的梯度传播和参数梯度累加方式。

至此，我们就可以把前面学习的卷积神经网络从一个统一的角度串联起来：卷积通过局部连接利用图像的局部空间结构，通过权值共享使相同的局部模式能够在不同位置重复检测，并因此具有平移等变性；池化进一步减弱特征对局部精确位置的敏感程度；而在训练过程中，这种特殊的网络结构又直接决定了梯度计算的形式。也正因如此，前面卷积层反向传播中那些看似复杂的逐位置累加，并不是额外增加的一套规则，而是卷积结构在矩阵梯度计算下的自然结果。到这里，卷积的结构、前向计算、空间性质以及反向传播就形成了一个完整的闭环。

\begin{quote}
    至此，我们已经掌握了卷积神经网络的基本原理。接下来，我们将从AlexNet开始，看看它如何引爆深度学习革命，以及卷积神经网络如何沿着更深、更强的方向不断发展。
\end{quote}